\conclusion

\section*{Основные результаты}

Решены все поставленные задачи и достигнута основная цель~--- разработан прототип системы автоматического мульти-задачного анализа музыкальных композиций.

\begin{enumerate}
  \item Реализованы и обучены пять архитектур мульти-задачных моделей (CNN с SE-блоками, CNN\,+\,BiLSTM, PANNs\,+\,Linear, PANNs\,+\,BiLSTM, AST) и проведены четыре последовательных эксперимента (v1, v2, AST~v2, CNN~v3 как negative result).
  \item Финальная модель AST~v2 показала лучший совокупный результат: среднее accuracy 58{,}5\,\%, сегмент 40{,}6\,\% (+5{,}3\,п.п. к CNN v2), жанр 81{,}5\,\%.
  \item Реализован композитный post-processing pipeline (Viterbi + musicological position priors + per-class min-duration), повышающий Boundary~F1 на 14\,пунктов и сокращающий число сегментов с 28 до 7{,}5 в среднем.
  \item Разработан веб-прототип с трёхзвенной архитектурой (FastAPI / Elysia\,+\,Bun / React\,+\,Vite), синхронизированным дашбордом из шести панелей и полноэкранным Vibe~Mode (WebGL2-шейдер).
  \item Идентифицированы ограничения: OOD-эффект на современной музыке, потолок per-window ($\sim 55$\,\%), размер AST (574\,МБ), отсутствие multi-label жанров.
\end{enumerate}

\section*{Направления дальнейшей работы}

\begin{enumerate}
  \item Temporal-декодер BiLSTM/CRF поверх AST-эмбеддингов~--- sequence-level предсказание структуры вместо per-window; по литературе даёт $+10$--$15$\,п.п. Boundary~F1.
  \item Дообучение жанровой головы на FMA-medium / MagnaTagATune с multi-label-предсказаниями~--- устраняет OOD-проблему.
  \item Knowledge distillation AST\,$\to$\,MobileViT для уменьшения модели в 10--20 раз и mobile deployment.
  \item Streaming inference через скользящее окно для треков произвольной длительности с фиксированной памятью.
  \item Real-time оптимизация ($<100$\,мс на окно) для DJ-инструментов и ритм-игр.
  \item Перспективно: joint training с генеративными задачами, self-supervised pretraining на музыкальных данных, cross-modal extension.
\end{enumerate}
